% =====================================================================
\chapter{El modo directo y los números duales}
\label{tema:2}
% =====================================================================

El modo directo propaga, junto con cada valor de la traza, su derivada respecto a una dirección fijada de las entradas. Este tema lo formula como propagación de tangentes, demuestra que calcula un producto jacobiano-vector y lo reinterpreta con los números duales, que reducen la regla de la cadena a aritmética. Se siguen la §3.1 de~\cite{baydin}, incluida la §3.1.1 sobre números duales, el capítulo~3 de~\cite{gw}, las notas de la clase~8 de MIT 18.337~\cite{mit18337} y la clase~5 de MIT 18.S096~\cite{mit18s096}.

\section{Propagación de tangentes}

Se fija una traza de $f\colon\R^{n}\to\R^{m}$ (\cref{def:traza}), un punto $x$ y una dirección $\dot x\in\R^{n}$.

\begin{definicion}[Tangentes]\label{def:tangente}
Las \emph{tangentes} de la traza son los números definidos por recurrencia como
\[
\dot v_{i-n}=\dot x_i\quad(i=1,\dots,n),\qquad
\dot v_i=\sum_{j\pre i}\frac{\partial\varphi_i}{\partial v_j}\,\dot v_j\quad(i=1,\dots,\ell),
\]
donde las derivadas parciales se evalúan en los valores de la traza. El vector $\dot x$ se llama \emph{semilla}.
\end{definicion}

\begin{teorema}[El modo directo calcula $J\dot x$]\label{teo:directo}
Sea $\gamma(t)=x+t\dot x$. Para cada $i$, $\dot v_i=\frac{d}{dt}v_i(\gamma(t))\big|_{t=0}=\nabla v_i(x)\T\dot x$, donde $v_i(\cdot)$ denota la variable $i$ como función de las entradas. En particular, el vector de tangentes de las salidas es $\dot y=J_f(x)\,\dot x$.
\end{teorema}

\addtocounter{definicion}{1}%

\begin{corolario}[Coste del jacobiano en modo directo]\label{cor:coste-directo}
Una pasada en modo directo da una columna de $J_f(x)$, la correspondiente a la semilla $e_k$. El jacobiano completo de $f\colon\R^{n}\to\R^{m}$ requiere $n$ pasadas, independientemente de $m$.
\end{corolario}

El corolario fija el campo de aplicación del modo directo: es óptimo cuando hay pocas entradas y muchas salidas y es inviable para el gradiente de una pérdida escalar respecto a $10^{11}$ parámetros, que exigiría $10^{11}$ pasadas.

\section{Números duales}

\begin{definicion}[Números duales]\label{def:duales}
El conjunto de los \emph{números duales} es $\Du=\{a+b\varepsilon : a,b\in\R\}$, con la suma y el producto
\[
(a+b\varepsilon)+(c+d\varepsilon)=(a+c)+(b+d)\varepsilon,\qquad
(a+b\varepsilon)(c+d\varepsilon)=ac+(ad+bc)\varepsilon .
\]
El número $a$ es la \emph{parte real} y $b$, la \emph{parte dual}.
\end{definicion}

El producto es el que resulta de multiplicar formalmente e imponer $\varepsilon^{2}=0$ con $\varepsilon\ne 0$. Con estas operaciones $\Du$ es un anillo conmutativo con unidad, isomorfo a $\R[\varepsilon]/(\varepsilon^{2})$, pero no es un cuerpo: $\varepsilon\cdot\varepsilon=0$ con $\varepsilon\ne 0$, de modo que $\varepsilon$ no tiene inverso~\cite[§3.1.1]{baydin}.

\begin{proposicion}[Inverso de un dual]\label{prop:inverso}
Si $a\ne 0$, el dual $a+b\varepsilon$ es invertible y
\[
\frac{1}{a+b\varepsilon}=\frac1a-\frac{b}{a^{2}}\,\varepsilon .
\]
\end{proposicion}

La parte dual del inverso es $b$ veces la derivada de $u\mapsto 1/u$ en $a$, y la parte dual del producto es la regla de Leibniz. Esto no es casual, y motiva la definición siguiente.

\begin{definicion}[Extensión dual de una operación elemental]\label{def:extension}
Sea $\varphi\colon U\subseteq\R^{k}\to\R$ de clase $C^{1}$. Su \emph{extensión dual} es
\[
\varphi(a_1+b_1\varepsilon,\dots,a_k+b_k\varepsilon)=\varphi(a)+\Bigl(\sum_{j=1}^{k}\frac{\partial\varphi}{\partial u_j}(a)\,b_j\Bigr)\varepsilon ,\qquad a\in U.
\]
Por ejemplo, $\exp(a+b\varepsilon)=e^{a}+e^{a}b\,\varepsilon$ y $\sen(a+b\varepsilon)=\sen a+(\cos a)\,b\,\varepsilon$.
\end{definicion}

\begin{proposicion}[Coherencia con la aritmética]\label{prop:coherencia}
Para la suma, el producto y la inversión, la extensión dual de la \cref{def:extension} coincide con las operaciones de la \cref{def:duales} y de la \cref{prop:inverso}.
\end{proposicion}

\begin{teorema}[Los duales propagan la derivada]\label{teo:duales}
Sea $f\colon\R^{n}\to\R^{m}$ dada por una traza y evalúese la traza en aritmética dual, con entradas $x_k+\dot x_k\varepsilon$ y la extensión dual de cada operación elemental. Entonces cada variable vale $v_i+\dot v_i\varepsilon$, con $\dot v_i$ la tangente de la \cref{def:tangente}, y las salidas valen
\[
f(x+\dot x\,\varepsilon)=f(x)+J_f(x)\,\dot x\,\varepsilon .
\]
\end{teorema}

El teorema justifica la lectura del desarrollo de Taylor del módulo~4 con un infinitesimal: $f(a+b\varepsilon)=f(a)+f'(a)b\varepsilon$ es el desarrollo de primer orden, y aquí el truncamiento es exacto porque los términos de orden superior contienen $\varepsilon^{2}=0$. Para evaluar $f'(a)$ basta con calcular $f(a+\varepsilon)$ y leer la parte dual.

\addtocounter{definicion}{1}%

\begin{observacion}[Duales con varias componentes]\label{obs:multidual}
Si se toman $p$ infinitesimales $\varepsilon_1,\dots,\varepsilon_p$ con $\varepsilon_r\varepsilon_s=0$ para todo $r,s$, cada variable lleva un vector de $p$ partes duales, y con las semillas $x_k+\sum_r(\dot x^{(r)})_k\varepsilon_r$ una sola pasada calcula $p$ productos $J\dot x^{(r)}$. Con $p=n$ y semillas $e_1,\dots,e_n$ se obtiene el jacobiano entero en una pasada vectorial, cuyo coste sigue siendo proporcional a $n\cdot\ops(f)$~\cite{mit18337}. La biblioteca JAX ofrece esta operación como \texttt{jvp} (\emph{Jacobian-vector product}).
\end{observacion}

\section{Coste del modo directo}

\begin{proposicion}\label{prop:coste-dual}
Existe una constante $\omega_{\mathrm{dir}}$, que depende solo del conjunto $\Phi$ de operaciones elementales, tal que una pasada en modo directo cuesta a lo sumo $\omega_{\mathrm{dir}}\cdot\ops(f)$ operaciones.
\end{proposicion}

\addtocounter{definicion}{1}%

